% Generated draft: requires mathematical and source review.
\documentclass[11pt]{article}
\usepackage[margin=25mm]{geometry}
\usepackage{fontspec}
\setmainfont{Latin Modern Roman}
\usepackage{amsmath,amssymb,mathtools,unicode-math}
\setmathfont{Latin Modern Math}
\usepackage{enumitem,needspace,xurl,hyperref}
\setlength{\parindent}{0pt}
\setlength{\parskip}{5pt}
\setlength{\emergencystretch}{3em}
\newcommand{\N}{\mathbb{N}}
\newcommand{\Z}{\mathbb{Z}}
\newcommand{\Q}{\mathbb{Q}}
\newcommand{\R}{\mathbb{R}}
\newcommand{\C}{\mathbb{C}}
\newcommand{\F}{\mathbb{F}}
\newcommand{\A}{\mathbb{A}}
\newcommand{\PP}{\mathbb{P}}
\newcommand{\E}{\mathbb{E}}
\newcommand{\eps}{\varepsilon}
\newcommand{\dd}{\,\mathrm{d}}
\begin{document}
\section*{UnsolvedMath ID: 10400026}
\textbf{Source problem code:} catalogue entry
\subsection{Definitions and mathematical statement}
Let \(X=\downarrow\!\downarrow\) be the source's two oriented interval strands and let \(A(X;\Z)\), \(A(S^1;\Z)\) be its integral diagram modules. They consist of integer linear combinations of diagrams on the indicated skeletons modulo the specified diagram relations. The exact chord/Jacobi, framing, and relation conventions are not included here and must not be changed when testing torsion.

Write \(\tau\in A(X;\Z)\) for the Dogolazky--Kneissler element shown in the source's Figure 7, reported to satisfy
\[
\tau\ne0,\qquad2\tau=0.
\]
For each admissible diagram embedding/insertion \(e\) into a circle-skeleton diagram, let \(j_e\) denote the induced integral-linear operation, and put \(x_e=j_e(\tau)\in A(S^1;\Z)\). Since the operation respects the relations,
\[
2x_e=j_e(2\tau)=0.
\]
Stanford's precise nonvanishing target is
\[
\exists e\text{ admissible},\qquad j_e(\tau)\ne0.
\]
Such an image would have exact order two. Nonzero source torsion does not itself imply a nonzero image, since the insertion/closure operation need not be injective.

The source says a positive answer would refute Conjecture 2.3, described in the supplied summary as torsion-freeness of \(A(S^1;\Z)\). That condition is
\[
\forall x\in A(S^1;\Z)\ \forall m\ge1,\qquad mx=0\Longrightarrow x=0.
\]
Conversely, vanishing of all images of this one element would not alone establish torsion-freeness for every circle diagram.

Figure 7's actual integral combination and the complete set of allowed embeddings/auxiliary diagram insertions are absent, so no concrete image or exhaustive test can be reconstructed. The trailing text ``2.2 Do Vassiliev invariants distinguish knots?'' is an adjacent section heading, not another clause of Question 2.4. The supplied research summary reports no known nontrivial image, without a verified exhaustive result.
\subsection{Short English statement}
Can some allowed embedding of the specified two-strand order-two diagram element give a nonzero circle-diagram class? Its image is automatically killed by two, but nonvanishing needs the missing Figure 7 and exact integral diagram conventions.
\subsection{Sources}
\begin{itemize}
\item Local UnsolvedMath registry, permanent ID \(10400026\), exact source code \texttt{catalogue entry}, original Question 2.4, attributed to T. Stanford.
\item Ohtsuki, \emph{Problems on invariants of knots and 3-manifolds} (2002), Question 2.4, recorded PDF page 29; \url{https://msp.org/gtm/2002/04/gtm-2002-04-024s.pdf}. Recorded access 2026-07-29; pdf-font-structure extraction.
\item Source Figure 7 for the Dogolazky--Kneissler element and Conjecture 2.3 for the circle-module torsion claim; the figure and full convention definitions are not supplied.
\item Supplied OPEN-TRIAGE summary reports no known nontrivial embedded image, with no additional precise publication citation.
\end{itemize}
\end{document}
